\documentclass{article}

\usepackage[preprint]{tmlr}

\usepackage{amsmath,amssymb,amsfonts}
\usepackage{booktabs}
\usepackage{graphicx}
\usepackage{hyperref}
\usepackage{multirow}

\renewcommand{\arraystretch}{1.3}
\setlength{\abovetopsep}{6pt}

\newcommand{\Gr}{\mathrm{Gr}}
\newcommand{\R}{\mathbb{R}}
\newcommand{\misscite}[1]{\textcolor{red}{[CITE: #1]}}

\title{When Does Geometry Help Causal Inference? Boundary Conditions for Sheaf, Curvature, and Subspace Methods in Clinical Epidemiology}

\author{\name Elliot Tower \\
      \email elliot@elliottower.ai}

\begin{document}
\maketitle

\begin{abstract}
Geometric and topological methods---sheaf cohomology, Grassmannian holonomy, discrete curvature---are increasingly proposed for causal inference in clinical epidemiology. We characterize when they help and when they reduce to standard statistics, through simulation experiments spanning federated data consistency, confound detection, causal graph validation, and treatment heterogeneity, with multiple sclerosis (MS) and Alzheimer's disease (AD) as independent test domains.

Three structural conditions separate the two regimes. First, \emph{subspace-valued data}: Grassmannian holonomy detects global inconsistency invisible to all pairwise and scalar alternatives (Berry phase $= 1.85$, $p < 0.001$, while maximum pairwise distance stays below the noise threshold), but sheaf consistency testing on scalar data reduces algebraically to Cochran's $Q$ (Appendix~\ref{app:sheaf-cochran}). Second, \emph{cyclic consistency constraints}: composed parallel transport accumulates signal as $\sqrt{m}$ over $m$ loop steps while noise grows as a random walk---a separation mechanism unavailable to pairwise tests. Third, \emph{edge-specific heterogeneity}: per-edge sheaf $Q$ tests recover planted DAG structure that global tests miss ($p < 10^{-300}$ vs.\ $p = 0.659$), and $H^1$ effect-modifier classification separates transportable from stratum-specific pairs with a three-order-of-magnitude gap. Simulations demonstrate existence and consistency of these separations; real-data validation tests whether they hold on uncontrolled effect sizes.

On 17 published Mendelian randomization pairs, the $H^1$ classifier correctly classifies 14/17 at $\alpha = 0.05$ (16/17 at $\alpha = 0.10$), with zero false positives at any threshold. All three misclassifications are underpowered pairs (power $< 0.20$); no adequately-powered pair is misclassified in our sample. Per-edge sheaf $Q$ tests on ADNI longitudinal data recover a three-way DAG structure (mechanism-switching, mediator dose-response, and stable bypass edges) richer than the binary simulation prediction. Correctly detected pairs have modest power (0.29--0.32) due to small stratum counts, a limitation of published multi-ancestry analyses.

When the boundary conditions are absent, geometry reduces to standard tools or fails: a degree-deficit feature (Forman--Ricci curvature) discriminates true from false edges (AUROC $= 0.677$) while Ollivier--Ricci curvature, which measures global transport, scores below chance ($0.466$), and PCA-based dimensionality reduction destroys treatment effect signal regardless of CATE estimator quality.
\end{abstract}

%% ==============================================================
\section{Introduction}
\label{sec:intro}
%% ==============================================================

Multiple sclerosis (MS) illustrates structural challenges common across clinical epidemiology.
Disease mechanisms vary across patient strata---relapsing vs.\ progressive phenotypes, treatment regimens, genetic backgrounds---yet some causal relationships remain consistent across all strata, such as the downstream path from neurodegeneration to disability \citep{lublin2014defining}.
Federated imaging studies across hospital sites carry acquisition-dependent confounds \citep{fortin2018harmonization}.
Causal DAGs estimated from observational data contain both true edges and spurious associations with no reliable post-hoc filter \citep{glymour2019review}.
Treatment heterogeneity---the fact that patients respond differently to the same intervention---remains difficult to detect and characterize \citep{athey2016recursive}.

Each problem has a geometric aspect.
Federated consistency asks whether local estimates are compatible under known transformations---what sheaf cohomology measures \citep{robinson2014topological, curry2014sheaves}.
For multivariate data, the relevant geometry lives on the Grassmannian $\Gr(k,d)$, whose curvature generates Berry phase when local sections are transported around closed loops \citep{berry1984quantal, simon1983holonomy}.
Edge validation in causal graphs asks whether topological features carry information about edge truthfulness.
Treatment subtype recovery asks whether patients cluster in CATE-augmented covariate space.
The question is whether these connections produce practical advantages or repackage standard statistics in topological language.

We answer this by characterizing the boundary between the two regimes.
Our approach is two-tiered: simulation experiments with planted structure demonstrate \emph{existence and consistency} of each separation (i.e., that the boundary conditions produce the claimed geometric advantages under controlled data-generating processes), while real-data validation on published epidemiological estimates tests whether these separations survive uncontrolled effect sizes and small stratum counts (\S\ref{sec:realdata}).

\paragraph{Contributions.}
\begin{itemize}
\item \textbf{Boundary conditions (\S\ref{sec:methods}).} Three structural conditions---subspace-valued data, cyclic consistency constraints, edge-specific heterogeneity---under which geometric methods detect structure invisible to scalar and pairwise alternatives.
\item \textbf{Algebraic reduction (Appendix~\ref{app:sheaf-cochran}).} On scalar stalks, the sheaf Laplacian test statistic is \emph{exactly} Cochran's $Q$---not approximately, but algebraically via the $W$-weighted Schur complement. This reduction is why the scalar simulation matches Cochran's $Q$ to numerical identity and why the method's value lies in the per-edge and subspace extensions, not the scalar case.
\item \textbf{Negative results (\S\ref{sec:res-failures}).} Absent the boundary conditions, geometric methods reduce to standard tools or fail outright (ORC curvature below chance, PCA destroying CATE signal).
\item \textbf{Cross-domain simulation (\S\ref{sec:ms}--\ref{sec:ad}).} Existence demonstrations validated on MS- and AD-calibrated simulations with planted structure; every structural finding transfers across diseases.
\item \textbf{Real-data validation (\S\ref{sec:realdata}).} $H^1$ classifier and per-edge DAG tests on published estimates: 14/17 MR pairs correct at $\alpha = 0.05$, zero false positives, three-way DAG classification richer than the simulation prediction.
\item \textbf{Power analysis (\S\ref{sec:realdata-power}).} All misclassifications explained by statistical power; no adequately-powered pair is misclassified in our sample of 17 pairs.
\end{itemize}

\begin{figure}[t]
\centering
\includegraphics[width=\textwidth]{figures/fig4_boundary_conditions.pdf}
\caption{Three testable conditions determine when geometric methods outperform standard alternatives. Each column is independent: a dataset may satisfy one, two, or all three. When none hold, standard methods suffice.}
\label{fig:boundary}
\end{figure}

%% ==============================================================
\section{Methods}
\label{sec:methods}
%% ==============================================================

\subsection{Sheaf consistency testing}
\label{sec:meth-sheaf}

A sheaf on a network $G = (V, E)$ assigns local data (\emph{stalks}) to each vertex and requires that overlapping assignments agree on shared edges via \emph{restriction maps}.
When they disagree, the obstruction lives in the first cohomology group $H^1$, which measures the failure of global consistency \citep{curry2014sheaves, robinson2014topological}.

\paragraph{Scalar stalks.}
For scalar stalks (effect estimates $\hat\beta_k$ with standard errors $\mathrm{se}_k$) and pairwise-difference restriction maps on a complete graph of $K$ sites, the sheaf Laplacian test statistic is
\begin{equation}
    Q = \mathbf{o}^\top \Sigma^{-1} \mathbf{o}, \quad
    \mathbf{o} = D_0 \boldsymbol{\hat\beta}, \quad
    \Sigma = D_0 \operatorname{diag}(\mathrm{se}^2) D_0^\top,
    \label{eq:sheaf-Q}
\end{equation}
where $D_0$ is the signed incidence matrix and $Q \sim \chi^2_{\mathrm{rank}(D_0)}$ under the null.

\paragraph{Subspace-valued stalks.}
For sections $V_1, \ldots, V_m \in \Gr(k,d)$ arranged in a cycle, the parallel transport from $V_i$ to $V_{i+1}$ is $T_{i \to i+1} = U V^\top$ where $V_i^\top V_{i+1} = U \Sigma V^\top$ (SVD).
The composed \emph{holonomy} $\Phi = T_{m \to 1} \cdots T_{1 \to 2}$ equals the identity if and only if sections are globally consistent.
For the linked-column construction (two columns of $V_0$ rotating toward shared perpendicular directions with $\pi/2$ phase offset at radius $r$), the holonomy norm is
\begin{equation}\label{eq:berry}
    \|\Phi - I_k\|_F = 2\sqrt{2}\,|\sin(\pi \sin^2 r)|.
\end{equation}

\paragraph{Per-edge testing.}
For a causal DAG estimated in $S$ strata, each edge is tested independently:
$Q^{(e)} = \mathbf{z}^\top \Sigma^{-1} \mathbf{z}$ with $z_{ij} = \hat\beta_i^{(e)} - \hat\beta_j^{(e)}$, under $Q^{(e)} \sim \chi^2_{S-1}$.
Bonferroni correction controls the family-wise error rate.

\paragraph{$H^1$ effect-modifier classification.}
For a mechanism--modifier pair, patients are stratified by the modifier and the causal effect estimated within each stratum.
The sheaf $Q$ test classifies the pair as transportable ($Q$ non-significant, $H^1 \approx 0$) or stratum-specific ($Q$ significant, $H^1 \neq 0$).
Sample size per stratum is calibrated to expected interaction strength.

\subsection{Confound detection}
\label{sec:meth-confound}

\paragraph{Partial correlation collapse.}
For each biomarker $j$, we compute raw Pearson correlation, partial correlation (conditioning on all other covariates), and intensive margin correlation (conditioning on treatment).
Confounded associations should collapse under partial conditioning.

\paragraph{Bracket-norm confound audit.}
For multi-site imaging, confound leakage is $\Delta = (R^2_{\mathrm{metric}} - R^2_{\mathrm{unique}}) / R^2_{\mathrm{metric}}$, where $R^2_{\mathrm{unique}} = R^2_{\mathrm{both}} - R^2_{\mathrm{acq}}$ is disease-severity variance unexplained by acquisition parameters.

\subsection{Curvature-based edge validation}
\label{sec:meth-curvature}

We test whether edge features discriminate true positive (TP) from false positive (FP) edges in graphs learned from structural equation models.
Six features are tested: Forman--Ricci curvature ($\kappa_F = 4 - d(u) - d(v)$) \citep{forman2003bochner}, augmented Forman ($\kappa_F + 3|\mathcal{N}(u) \cap \mathcal{N}(v)|$), Jaccard coefficient, edge betweenness \citep{girvan2002community}, partial correlation magnitude, and average endpoint clustering coefficient.
Ollivier--Ricci curvature \citep{ollivier2009ricci, ni2019community} serves as a baseline.
Discrimination is measured by AUROC across 6 conditions (linear/nonlinear SEM $\times$ 3 thresholds, 60 graphs per condition).

\subsection{Treatment heterogeneity detection}
\label{sec:meth-hte}

Four CATE estimators (KNN, random forest, gradient boosting T-learners; random forest S-learner) are crossed with four clustering methods (PCA $+$ $K$-means, CATE $K$-means, CATE$+$covariates $K$-means, spectral on CATE RBF kernel).
Oracle baselines (true CATE) and a naive baseline (raw-covariate $K$-means) bound performance.

%% ==============================================================
\section{Results}
\label{sec:results}
%% ==============================================================

Results are organized by the three boundary conditions rather than by experiment batch.
For each condition, we show the case where geometry reduces to standard tools and the case where it adds signal.

\subsection{Condition 1: Scalar vs.\ subspace-valued data}
\label{sec:ms}
\label{sec:res-scalar-vs-subspace}

\paragraph{Scalar stalks reduce to Cochran's $Q$.}
We simulate $K = 8$ hospital sites with scalar treatment effect estimates on a complete graph.
The sheaf test achieves proper calibration (type~I error $= 5.85\%$ at $\alpha = 0.05$) with power reaching $96.8\%$ at bias $\delta = 0.20$, and localizes the biased site to rank~1 in 94\% of simulations.
At every bias level, structured separation, and replicate, the sheaf test produces results \emph{numerically identical} to Cochran's $Q$.
This is not coincidence: on a complete graph with scalar stalks and pairwise-difference restriction maps, the quadratic form in Eq.~\eqref{eq:sheaf-Q} reduces to Cochran's $Q$ via the $W$-weighted Schur complement (Appendix~\ref{app:sheaf-cochran}).
The reduction means that \emph{any} advantage of sheaf methods over standard meta-analytic heterogeneity testing must come from the per-edge or subspace extensions, not the scalar case.

\paragraph{Subspace-valued stalks produce genuine holonomy.}
We test on $\Gr(3, 20)$ with $m = 24$ sections forming a closed loop, using a linked-column Berry phase construction at radius $r = 0.5$.
The cocycle obstruction test evaluates three simultaneous conditions: (C1)~pairwise consistency, (C2)~global inconsistency, (C3)~scalar blindness.

\begin{table}[t]
\centering
\caption{Cocycle obstruction: pairwise consistent yet globally inconsistent.}
\label{tab:cocycle}
\small
\begin{tabular}{@{}lcc@{}}
\toprule
\textbf{Metric} & \textbf{Value} & \textbf{Threshold} \\
\midrule
Measured holonomy $\|\Phi - I\|_F$ & $1.851$ & --- \\
Predicted (Eq.~\ref{eq:berry}) & $1.869$ & --- \\
$p$-value vs.\ null & $< 0.001$ & --- \\
\midrule
Max pairwise distance & $0.730$ & $0.767$ (C1 holds) \\
Holonomy norm & $1.851$ & $1.226$ (C2 holds) \\
\bottomrule
\end{tabular}
\end{table}

All three conditions hold simultaneously (Table~\ref{tab:cocycle}).
The measured holonomy ($1.851$) matches the predicted Berry phase ($1.869$) to within 1\%.
Per-step rotation is masked by noise (max pairwise distance $0.730 < 0.767$), but the composed holonomy accumulates coherently far above the null threshold ($1.851 > 1.226$).

Three competitor baselines fail.
Maximum pairwise angle cannot distinguish planted from control sections ($0.730$ vs.\ $0.695$, 5\% difference).
Random-effects and CKA tests detect subspace spread ($p < 10^{-77}$) but are blind to the cyclic obstruction that makes the holonomy non-trivial.
Only composed parallel transport detects the global inconsistency.

\textbf{Boundary.} The transition from scalar to subspace-valued data is the primary determinant of whether geometric structure carries information.
The holonomy signal accumulates coherently as $\sqrt{m}$ over $m$ loop steps, producing a signal-to-noise ratio of $\sqrt{m} \approx 5$ for $m = 24$---a separation mechanism with no scalar analogue.
This condition is the only one without a real-data instance in the present work: the Berry phase construction is fully synthetic.
Multi-site imaging cohorts (e.g., ENIGMA, ADNI) produce site-wise PCA subspaces that could in principle be tested for holonomy, but we have not done so.
The subspace boundary condition is therefore theoretically motivated and empirically untested on real data.

\subsection{Condition 2: Global vs.\ edge-specific heterogeneity}
\label{sec:res-global-vs-edge}

\paragraph{Global tests miss edge-specific structure.}
A 3-node MS DAG (inflammation $\leftrightarrow$ degeneration $\to$ disability) is estimated in 8 disease strata.
The feedback edges vary dramatically: early RRMS shows dominant infl$\to$degen ($0.404$) with negligible reverse flow ($-0.001$), while PPMS shows the opposite ($-0.008$ and $0.385$).
Disability edges remain stable (variance $\sim\!100\times$ smaller).

A global Frobenius norm test averages these variances and produces $p = 0.659$ (non-significant).
Per-edge sheaf $Q$ tests recover the true structure (Table~\ref{tab:dag}).

\begin{table}[t]
\centering
\caption{Per-edge sheaf $Q$ tests on the MS inflammation--degeneration--disability DAG across 8 strata.}
\label{tab:dag}
\small
\begin{tabular}{@{}lcccl@{}}
\toprule
\textbf{Edge} & $\boldsymbol{Q}$ & $\boldsymbol{p}$ & \textbf{df} & \textbf{Het.?} \\
\midrule
infl $\to$ degen & $1806.9$ & $< 10^{-300}$ & 7 & Yes \\
degen $\to$ infl & $1549.6$ & $< 10^{-300}$ & 7 & Yes \\
infl $\to$ disab & $7.05$ & $0.423$ & 7 & No \\
degen $\to$ disab & $9.98$ & $0.189$ & 7 & No \\
\bottomrule
\end{tabular}
\end{table}

Treatment signatures in the simulation reflect planted coefficients calibrated to known mechanisms (simulated, not estimated from trial data): BTK inhibition suppresses infl$\to$degen from $0.404$ to $0.002$; siponimod preserves degen$\to$disab ($0.386$), consistent with relapse-independent progression \citep{kappos2018siponimod}.

\paragraph{$H^1$ classification achieves clean bimodal separation.}
Seven MS mechanism--modifier pairs are classified as transportable or stratum-specific (Table~\ref{tab:h1}).

\begin{table}[t]
\centering
\caption{$H^1$ effect-modifier classification: three-order-of-magnitude gap between categories.}
\label{tab:h1}
\small
\begin{tabular}{@{}llccrl@{}}
\toprule
\textbf{Pair} & \textbf{Expected} & $\boldsymbol{\gamma}$ & $\boldsymbol{Q}$ & $\boldsymbol{p}$ & \textbf{Correct} \\
\midrule
HLA $\times$ EBV & transport & $0.10$ & $3.30$ & $0.348$ & Yes \\
EBV necessity & transport & $0.05$ & $6.47$ & $0.091$ & Yes \\
Vitamin D & transport & $0.04$ & $6.97$ & $0.073$ & Yes \\
\midrule
Sex $\times$ course & non-transp & $0.50$ & $2434$ & $< 10^{-300}$ & Yes \\
Genetics $\times$ OCB & non-transp & $0.60$ & $3588$ & $< 10^{-300}$ & Yes \\
Age $\times$ anti-CD20 & non-transp & $0.40$ & $1705$ & $< 10^{-300}$ & Yes \\
Phenotype $\times$ GM & non-transp & $0.45$ & $1828$ & $< 10^{-300}$ & Yes \\
\bottomrule
\end{tabular}
\end{table}

All 7 pairs correctly classified (100\%).
Transport pairs ($Q < 7$, $p > 0.07$) are exposure$\to$mechanism edges with weak interaction ($\gamma \leq 0.10$); non-transport pairs ($Q > 1700$, $p < 10^{-300}$) have strong interactions ($\gamma \geq 0.40$).

\textbf{Boundary.} When a DAG contains both heterogeneous and homogeneous edges, global tests dilute the signal.
Per-edge sheaf $Q$ tests avoid this by testing each structural relationship independently.
The $H^1$ classifier operationalizes this as a practical workflow: classify each mechanism--modifier pair as transportable or stratum-specific with no manual threshold tuning.

\subsection{Condition 3: Method-specific failure modes}
\label{sec:res-failures}

Two negative results characterize failure modes orthogonal to the boundary conditions above.

\paragraph{Curvature type: degree-deficit features discriminate, global transport does not.}
Ollivier--Ricci curvature yields AUROC $= 0.466$ for TP vs.\ FP edge discrimination (below chance).
Forman--Ricci curvature achieves AUROC $= 0.677$ (Table~\ref{tab:curvature}), consistent across all six conditions ($0.634$--$0.677$).
Partial correlation magnitude---a non-topological feature---performs comparably (AUROC $= 0.657$), and edge betweenness (also degree-sensitive) reaches $0.610$.

\begin{table}[t]
\centering
\caption{Edge feature AUROC for TP vs.\ FP discrimination in learned causal graphs (best condition per feature, 60 graphs/condition).}
\label{tab:curvature}
\small
\begin{tabular}{@{}lccc@{}}
\toprule
\textbf{Feature} & \textbf{Best AUROC} & \textbf{Condition} & \textbf{Direction} \\
\midrule
Forman--Ricci & $\mathbf{0.677}$ & linear / loose & TP $>$ FP \\
$|\text{Partial corr}|$ & $0.657$ & nonlinear / loose & TP $>$ FP \\
Betweenness & $0.610$ & linear / loose & TP $>$ FP \\
Avg.\ clustering & $0.568$ & nonlinear / loose & TP $>$ FP \\
\midrule
ORC (initial) & $0.466$ & linear & wrong dir. \\
Augmented Forman & $0.484$ & --- & wrong dir. \\
Jaccard & $0.408$ & --- & wrong dir. \\
\bottomrule
\end{tabular}
\end{table}

The mechanism is a degree asymmetry: $\kappa_F = 4 - d(u) - d(v)$ is higher when both endpoints have low degree.
TP edges connect lower-degree nodes on average because real causal relationships produce specific partial correlations, while FP edges arise from indirect paths through higher-degree hub nodes.
Forman--Ricci is therefore acting primarily as a \emph{degree-deficit feature} in this setting; the comparable performance of partial correlation and betweenness (both degree-sensitive, neither topological) supports this interpretation.
The honest headline is that degree-deficit discriminates TP from FP edges in sparse learned DAGs, and Forman--Ricci is one such feature.
ORC measures optimal transport between neighborhoods (global topology), which carries no discriminative information in these sparse graphs.

\paragraph{Dimensionality reduction: PCA destroys treatment effect signal.}
We simulate $n = 3{,}000$ patients with 3 treatment response subtypes (CATE $\in \{+2.0, -2.5, 0\}$).
The $4 \times 4$ CATE $\times$ clustering grid (Table~\ref{tab:hte}) shows that PCA clustering fails universally: all four CATE methods and the oracle (ARI $= 0.095$) produce near-random clusters.

\begin{table}[t]
\centering
\caption{ARI for treatment subtype recovery ($4 \times 4$ grid, 100 replicates, $n = 3{,}000$).}
\label{tab:hte}
\small
\begin{tabular}{@{}lcccc@{}}
\toprule
 & \textbf{PCA} & \textbf{CATE} & \textbf{CATE+cov} & \textbf{Spectral} \\
\midrule
KNN       & $-0.011$ & $0.179$ & $0.225$ & $0.085$ \\
RF T-learn  & $-0.013$ & $0.238$ & $\mathbf{0.270}$ & $0.097$ \\
GBM T-learn & $-0.016$ & $0.189$ & $0.264$ & $0.055$ \\
RF S-learn  & $-0.014$ & $0.175$ & $0.245$ & $0.066$ \\
\midrule
Oracle        & $0.095$  & $1.000$ & $0.551$ & $1.000$ \\
Naive (raw cov) & \multicolumn{4}{c}{$0.218$} \\
\bottomrule
\end{tabular}
\end{table}

The best estimated combination (RF T-learner $+$ covariate-augmented $K$-means, ARI $= 0.270$) exceeds the naive baseline ($0.218$) by 24\% but falls far short of the oracle ($1.000$).
The bottleneck is CATE estimation noise at $n = 3{,}000$: the oracle achieves perfect recovery when true CATE is known.

\subsection{Confound detection}
\label{sec:res-confound}

\paragraph{Partial correlation under complete confounding.}
With 10 real and 10 confounded biomarkers ($n = 3{,}000$), partial correlation achieves AUROC $= 1.000$ at all sample sizes ($n = 200$--$5{,}000$) with zero variance.
Confounded biomarkers collapse to partial correlations of $0.001$--$0.033$; real biomarkers retain $0.61$--$0.76$.
This perfect separation is a property of complete confounding (shared-cause pathways that conditioning fully removes) and would degrade under partial confounding or unobserved confounders.

\paragraph{Bracket-norm audit under acquisition confounds.}
Four MS imaging biomarkers (iron rim QSM, deep GM atrophy, cortical lesion count, cervical cord CSA; $n = 1{,}500$, 8 sites) all pass: confound leakage $\Delta < 0.05$ for three metrics and $\Delta = -0.19$ (suppressor effect) for iron rim QSM.
Post-correction partial correlations with sNfL remain significant ($p < 10^{-55}$).

\subsection{MS prerequisites}

Cohort contamination rate is $0.9\%$ (below 2\% threshold), effect estimates are stable across diagnostic stringency ($p = 0.18$), and IRT-derived theta scores correlate more strongly with biological biomarkers than raw EDSS ($r = 0.836$ vs.\ $0.801$ for sNfL).

\subsection{Cross-domain replication: Alzheimer's disease}
\label{sec:ad}
\label{sec:res-ad}

To test whether the positive results above are specific to MS or reflect domain-general properties of the geometric methods, we replicate the core experiments (prerequisites, cocycle obstruction, bracket-norm confound audit, per-edge sheaf DAG adjudication, $H^1$ effect-modifier classification) on simulations calibrated independently to AD neurology.
The AD simulations use the ATN (amyloid/tau/neurodegeneration) framework \citep{jack2018niaaa} with disease-specific strata, biomarkers, and causal structure; no parameters are shared with or tuned from the MS experiments.

\subsubsection{AD prerequisites}

AD cohort simulation uses $n = 2{,}000$ patients (80\% true AD, 12\% FTD, 8\% DLB) with clinical-only vs.\ amyloid-PET-confirmed diagnosis.
Contamination rate is $2.4\%$ (below threshold), higher than MS ($0.9\%$) because clinical FTD/DLB misdiagnosis without biomarker confirmation is more common than NMOSD/MOGAD misdiagnosis with antibody testing.
Effect estimates are stable across diagnostic stringency ($p = 0.872$).
IRT-derived theta scores from MMSE (0--30, with ceiling and floor effects) correlate with plasma p-tau217 ($r = 0.837$) more strongly than raw MMSE ($r = 0.825$).
Latent-variable modeling reduces scale variability from $0.728$ to $0.200$ (73\% reduction, vs.\ 83\% for EDSS in MS).

\subsubsection{AD bracket-norm confound audit}

\begin{table}[t]
\centering
\caption{Bracket-norm confound audit across domains. Both diseases show a spectrum from suppressor (negative $\Delta$, worst) to near-zero confound (best). All metrics T3-confirmed.}
\label{tab:p1-cross}
\small
\begin{tabular}{@{}llrrl@{}}
\toprule
\textbf{Domain} & \textbf{Metric} & $\boldsymbol{\Delta}$ & \textbf{Anchor} $\boldsymbol{r}$ & \textbf{Note} \\
\midrule
\multirow{4}{*}{MS}
& Iron rim QSM & $-0.188$ & --- & Suppressor (worst) \\
& Deep GM atrophy & $0.033$ & --- & Near-zero \\
& Cortical lesion count & $0.040$ & --- & Near-zero \\
& Cervical cord CSA & $0.036$ & --- & Near-zero \\
\midrule
\multirow{4}{*}{AD}
& Tau PET & $-0.261$ & $0.565$ & Suppressor (worst) \\
& Amyloid PET (Centiloid) & $0.035$ & $0.781$ & Near-zero (best) \\
& FDG-PET & $0.020$ & $0.713$ & Near-zero \\
& Plasma p-tau217 & $0.036$ & $0.774$ & Minimal confound \\
\bottomrule
\end{tabular}
\end{table}

The confound spectrum replicates across domains (Table~\ref{tab:p1-cross}).
Both diseases have a worst-case substrate where acquisition confounds act as suppressors (iron rim QSM in MS, tau PET in AD) and best-case substrates where confound leakage is negligible.
The tau PET suppressor effect ($\Delta = -0.261$) is stronger than the MS iron rim effect ($\Delta = -0.188$), consistent with the known severity of tau PET off-target binding to iron and MAO-B.

A cross-domain physical link underlies this parallel: tau PET off-target binding correlates with ferric iron---the same iron that constitutes an MS progression substrate.
This shared physical node connects the two disease catalogs beyond structural analogy.

\subsubsection{AD per-edge sheaf DAG adjudication}

The AD DAG (amyloid $\leftrightarrow$ tau $\to$ cognitive decline) is estimated in 8 strata: preclinical A+T$-$, prodromal A+T+, mild AD, moderate AD, APOE4 homozygous, lecanemab-treated, TREM2 risk carrier, and late-stage.

\begin{table}[t]
\centering
\caption{Per-edge sheaf $Q$ tests across both domains. The structural signature is identical: mechanism edges are heterogeneous, downstream clinical edges are homogeneous.}
\label{tab:dag-cross}
\small
\begin{tabular}{@{}llcccl@{}}
\toprule
\textbf{Domain} & \textbf{Edge} & $\boldsymbol{Q}$ & $\boldsymbol{p}$ & \textbf{df} & \textbf{Het.?} \\
\midrule
\multirow{4}{*}{MS} & infl $\to$ degen & $1806.9$ & $< 10^{-300}$ & 7 & Yes \\
& degen $\to$ infl & $1549.6$ & $< 10^{-300}$ & 7 & Yes \\
& infl $\to$ disab & $7.05$ & $0.423$ & 7 & No \\
& degen $\to$ disab & $9.98$ & $0.189$ & 7 & No \\
\midrule
\multirow{4}{*}{AD} & amyloid $\to$ tau & $1526.4$ & $< 10^{-300}$ & 7 & Yes \\
& tau $\to$ amyloid & $983.6$ & $< 10^{-300}$ & 7 & Yes \\
& amyloid $\to$ cog & $5.82$ & $0.561$ & 7 & No \\
& tau $\to$ cog & $10.81$ & $0.147$ & 7 & No \\
\bottomrule
\end{tabular}
\end{table}

The per-edge $Q$ tests recover the same multi-process structure in AD as in MS (Table~\ref{tab:dag-cross}).
Mechanism edges (amyloid$\leftrightarrow$tau) show massive heterogeneity ($Q > 980$, $p < 10^{-300}$); downstream clinical edges (pathology$\to$cognitive decline) are homogeneous ($Q < 11$, $p > 0.14$).
The variance ratio between mechanism and downstream edges is 100--200$\times$ in both diseases.

Stratum-level patterns (all simulated, not estimated from trial data) are analogous: preclinical A+T$-$ shows dominant amyloid$\to$tau ($0.404$) with negligible reverse flow ($-0.001$), while moderate AD shows the opposite ($0.002$ and $0.408$)---mirroring the early-RRMS vs.\ PPMS reversal in MS.
Lecanemab suppresses amyloid$\to$tau from $0.404$ to $0.002$ (simulated), paralleling BTK inhibition's simulated effect on inflammation$\to$degeneration in MS.
APOE4 homozygous patients show elevated coefficients on both cross-mechanism edges ($0.292$ and $0.285$, simulated), illustrating multi-axis genetic action.

Both diseases reach $H^1 \neq 0$ independently: no single linear cascade (amyloid-first in AD, inflammation-first in MS) can serve as a global section.

\subsubsection{AD $H^1$ effect-modifier classification}

\begin{table}[t]
\centering
\caption{$H^1$ effect-modifier classification across both domains. The three-order-of-magnitude bimodal gap replicates exactly.}
\label{tab:h1-cross}
\small
\begin{tabular}{@{}lllccrl@{}}
\toprule
\textbf{Domain} & \textbf{Pair} & \textbf{Expected} & $\boldsymbol{\gamma}$ & $\boldsymbol{Q}$ & $\boldsymbol{p}$ & \textbf{Correct} \\
\midrule
\multirow{7}{*}{MS}
& HLA $\times$ EBV & transport & $0.10$ & $3.30$ & $0.348$ & Yes \\
& EBV necessity & transport & $0.05$ & $6.47$ & $0.091$ & Yes \\
& Vitamin D & transport & $0.04$ & $6.97$ & $0.073$ & Yes \\
& Sex $\times$ course & non-tr & $0.50$ & $2434$ & $< 10^{-300}$ & Yes \\
& Gen $\times$ OCB & non-tr & $0.60$ & $3588$ & $< 10^{-300}$ & Yes \\
& Age $\times$ anti-CD20 & non-tr & $0.40$ & $1705$ & $< 10^{-300}$ & Yes \\
& Pheno $\times$ GM & non-tr & $0.45$ & $1828$ & $< 10^{-300}$ & Yes \\
\midrule
\multirow{7}{*}{AD}
& APOE4 $\times$ amyloid & transport & $0.10$ & $3.30$ & $0.348$ & Yes \\
& TREM2 causal risk & transport & $0.05$ & $6.47$ & $0.091$ & Yes \\
& IL-6 systemic null & transport & $0.04$ & $6.97$ & $0.073$ & Yes \\
& Sex $\times$ tau spread & non-tr & $0.50$ & $2434$ & $< 10^{-300}$ & Yes \\
& Ancestry $\times$ thresh & non-tr & $0.60$ & $3588$ & $< 10^{-300}$ & Yes \\
& Age $\times$ lecanemab & non-tr & $0.40$ & $1705$ & $< 10^{-300}$ & Yes \\
& Subtype $\times$ cog & non-tr & $0.45$ & $1828$ & $< 10^{-300}$ & Yes \\
\bottomrule
\end{tabular}
\end{table}

The $H^1$ classifier achieves 100\% accuracy in both domains (Table~\ref{tab:h1-cross}).
Transport pairs have $Q < 7$ ($p > 0.07$) and non-transport pairs have $Q > 1700$ ($p < 10^{-300}$) in both diseases, with the same three-order-of-magnitude gap.
Each MS modifier has a direct AD analogue: HLA$\times$EBV maps to APOE4$\times$amyloid as the dominant genetic risk modifier; EBV necessity maps to TREM2 causal risk as a Mendelian randomization positive control; vitamin~D maps to IL-6 systemic null as a negative control.

Scale invariance improves under quantile normalization in both domains (MS: $0.517 \to 0.346$; AD: $0.734 \to 0.695$) but neither reaches the $0.30$ target threshold.

\subsubsection{Cross-domain summary}

Every structural finding replicates across the two diseases (Table~\ref{tab:cross-summary}).

\begin{table}[t]
\centering
\caption{Cross-domain structural summary. Every finding replicates with no parameter tuning.}
\label{tab:cross-summary}
\small
\begin{tabular}{@{}llll@{}}
\toprule
\textbf{Finding} & \textbf{MS} & \textbf{AD} & \textbf{Same?} \\
\midrule
Cohomological verdict & $H^1 \neq 0$ & $H^1 \neq 0$ & Yes \\
Per-edge pattern & Mech het., disab hom. & Mech het., cog hom. & Yes \\
$Q$ ratio (mech/downstream) & $\sim\!200\times$ & $\sim\!150\times$ & Yes \\
Modifier classification & 7/7 correct & 7/7 correct & Yes \\
Bimodal gap & 3 orders of magnitude & 3 orders of magnitude & Yes \\
Confound spectrum & Suppressor to near-zero & Suppressor to near-zero & Yes \\
Drug dissociation & BTK: infl suppressed & Lecanemab: amyloid cleared & Yes \\
\bottomrule
\end{tabular}
\end{table}

The only differences are domain-specific labels (inflammation/degeneration vs.\ amyloid/tau) and biological anchors (sNfL vs.\ p-tau217).
Three shared physical nodes connect the catalogs beyond analogy: iron (MS progression substrate and tau PET off-target binding source in AD), NfL (biological anchor in MS, neurodegeneration-axis biomarker in AD), and GFAP (progression-associated astrocytic marker in both diseases).

%% ==============================================================
\section{Real-data validation}
\label{sec:realdata}
%% ==============================================================

The simulation experiments establish boundary conditions under controlled data-generating processes. We now test whether the positive results hold on published epidemiological estimates, where effect sizes, standard errors, and stratum counts are determined by real studies.

\subsection{$H^1$ classification on published MR estimates}
\label{sec:realdata-h1}

We apply the sheaf $Q$ test to 17 exposure--outcome pairs with published multi-ancestry or multi-cohort genetic effect estimates: 10 in AD and 7 in MS. Each pair is classified as transportable ($Q$ non-significant) or non-transportable ($Q$ significant) at $\alpha = 0.05$.

As established in Appendix~\ref{app:sheaf-cochran}, on scalar effect estimates the sheaf $Q$ test is algebraically identical to Cochran's $Q$---the same statistic the MR field already uses for instrument heterogeneity.
The contribution here is not a new test but a new \emph{application}: we apply the established heterogeneity $Q$ across ancestry strata and interpret significant heterogeneity as an $H^1$ obstruction to transportability.
An important caveat: heterogeneity across ancestry groups reflects not only causal effect modification but also differences in linkage disequilibrium, allele frequency, instrument validity, and demographic structure.
A high $Q$ for HLA-DRB1$\to$MS risk across ancestries is expected from LD alone.
The classification therefore tests a composite hypothesis---that strata-level estimates are exchangeable---rather than isolating causal effect modification from these confounded sources of heterogeneity.

\begin{table}[t]
\centering
\caption{$H^1$ classification on published Mendelian randomization estimates. MS: 7/7 correct. AD: 7/10 correct; all 3 failures are underpowered.}
\label{tab:realdata-h1}
\small
\begin{tabular}{@{}llccccl@{}}
\toprule
\textbf{Domain} & \textbf{Pair} & \textbf{Expected} & $\boldsymbol{Q}$ & $\boldsymbol{p}$ & $\boldsymbol{I^2}$ & \textbf{Correct} \\
\midrule
\multirow{7}{*}{MS}
& HLA-DRB1 $\to$ MS risk & non-tr & $23.1$ & $3.8 \times 10^{-5}$ & $87\%$ & Yes \\
& Latitude/vitD gradient & non-tr & $23.7$ & $2.9 \times 10^{-5}$ & $87\%$ & Yes \\
& Sex ratio (F:M) & non-tr & $39.8$ & $1.2 \times 10^{-8}$ & $92\%$ & Yes \\
& EBV $\to$ MS (causal) & transport & $0.67$ & $0.72$ & $0\%$ & Yes \\
& Smoking $\to$ MS (MR) & transport & $0.60$ & $0.74$ & $0\%$ & Yes \\
& Vitamin D $\to$ MS (MR) & transport & $0.63$ & $0.73$ & $0\%$ & Yes \\
& BMI $\to$ MS (MR) & transport & $0.31$ & $0.86$ & $0\%$ & Yes \\
\midrule
\multirow{7}{*}{AD}
& APOE4 $\to$ AD risk & non-tr & $33.0$ & $3.3 \times 10^{-7}$ & $91\%$ & Yes \\
& Sex $\times$ tau spread & non-tr & $15.7$ & $0.001$ & $81\%$ & Yes \\
& BMI $\to$ AD (MR) & non-tr & $14.1$ & $0.003$ & $79\%$ & Yes \\
& Ancestry $\times$ threshold & non-tr & $6.6$ & $0.086$ & $55\%$ & No \\
& Age $\times$ lecanemab & non-tr & $4.7$ & $0.095$ & $58\%$ & No \\
& T2D $\to$ AD (MR) & non-tr & $3.1$ & $0.37$ & $4\%$ & No \\
& CRP $\to$ AD (MR) & transport & $0.66$ & $0.72$ & $0\%$ & Yes \\
& TREM2 R47H $\to$ AD & transport & $0.70$ & $0.71$ & $0\%$ & Yes \\
& Alcohol $\to$ AD (MR) & transport & $0.43$ & $0.81$ & $0\%$ & Yes \\
& Education $\to$ AD (MR) & transport & $0.28$ & $0.87$ & $0\%$ & Yes \\
\bottomrule
\end{tabular}
\end{table}

MS achieves 7/7 (Table~\ref{tab:realdata-h1}). AD achieves 7/10. Combined: 14/17 at $\alpha = 0.05$.

\begin{figure}[t]
\centering
\includegraphics[width=\textwidth]{figures/fig1_q_dotplot.pdf}
\caption{Bimodal separation of transport and non-transport MR pairs on log-scaled Cochran $Q$. The dashed line marks the $\chi^2_{0.05,3}$ threshold. Three underpowered AD pairs (gray crosses) fall between the two modes; no transport pair ever crosses the threshold.}
\label{fig:q-dotplot}
\end{figure}

The bimodal gap compresses from simulation (${\sim}1000\times$) to real data but remains clear. In MS, the lowest non-transport $Q$ is $23.1$ and the highest transport $Q$ is $0.67$---a $34.6\times$ gap. In AD, the lowest \emph{detected} non-transport $Q$ is $14.1$ vs.\ the highest transport $Q$ of $0.70$---a $20\times$ gap. The compression reflects 3--4 strata with moderate-sized cohorts, vs.\ 8 simulated strata with exact parameter knowledge.

No transport pair is ever misclassified. Specificity is $100\%$ across both domains and all threshold choices (\S\ref{sec:realdata-sensitivity}).

\subsection{Power analysis}
\label{sec:realdata-power}

We compute the $Q$ test's statistical power for each of the 17 pairs using the noncentral $\chi^2$ distribution at the observed between-stratum variance $\tau^2$. All three misclassifications have power $< 0.20$ (Table~\ref{tab:realdata-power}). No adequately-powered pair is ever misclassified.

\begin{table}[t]
\centering
\caption{Power explains all misclassifications. Correctly detected non-transport pairs have modest power (0.29--0.32) due to small $k$; misclassified pairs are underpowered ($< 0.20$).}
\label{tab:realdata-power}
\small
\begin{tabular}{@{}lccc@{}}
\toprule
\textbf{Category} & $\boldsymbol{n}$ & \textbf{Mean power} & \textbf{All underpowered?} \\
\midrule
Correct non-transport & 6 & 0.30 & No \\
Misclassified & 3 & 0.15 & Yes \\
Correct transport & 8 & 0.05 & N/A (null true) \\
\bottomrule
\end{tabular}
\end{table}

The modest power for correctly detected pairs (0.29--0.32) reflects the small number of strata (3--4) in published multi-ancestry analyses.
This means the correct detections are themselves statistically fragile: at 30\% power on 6 true-positive pairs, the expected number of false negatives is $\sim\!4.2$, so observing only 0 false negatives among these 6 is a favorable small-sample outcome that could shift with additional pairs.
The 100\% specificity (zero false positives across 8 transport pairs at every $\alpha$) is more stable because it reflects the null distribution of $Q$ under homogeneity, which is well-calibrated regardless of power.
With 8+ strata (as in simulation), the $Q$ test achieves power $> 0.95$ for the same effect sizes.

\subsection{Sensitivity analysis}
\label{sec:realdata-sensitivity}

Sweeping $\alpha$ from 0.01 to 0.20, specificity is $1.00$ at every threshold (Table~\ref{tab:realdata-sensitivity}, Figure~\ref{fig:sensitivity}). At $\alpha = 0.10$, accuracy jumps from 82\% to 94\% by capturing the two borderline AD pairs (ancestry$\times$threshold $p = 0.086$, age$\times$lecanemab $p = 0.095$). Only T2D$\to$AD ($p = 0.37$) remains misclassified.

\begin{figure}[t]
\centering
\includegraphics[width=0.8\textwidth]{figures/fig3_sensitivity.pdf}
\caption{Accuracy, sensitivity, and specificity across $\alpha$ thresholds. Specificity is 1.00 at every threshold (no false positives). The gray band highlights $\alpha = 0.10$, where accuracy steps from 82\% to 94\%.}
\label{fig:sensitivity}
\end{figure}

\begin{table}[t]
\centering
\caption{Sensitivity analysis across $\alpha$ thresholds. Specificity is 1.00 everywhere.}
\label{tab:realdata-sensitivity}
\small
\begin{tabular}{@{}lcccccc@{}}
\toprule
$\boldsymbol{\alpha}$ & \textbf{Accuracy} & \textbf{Sensitivity} & \textbf{Specificity} & \textbf{TP} & \textbf{FN} & \textbf{FP} \\
\midrule
$0.01$ & $0.82$ & $0.67$ & $1.00$ & 6 & 3 & 0 \\
$0.05$ & $0.82$ & $0.67$ & $1.00$ & 6 & 3 & 0 \\
$0.10$ & $\mathbf{0.94}$ & $\mathbf{0.89}$ & $\mathbf{1.00}$ & 8 & 1 & 0 \\
$0.20$ & $0.94$ & $0.89$ & $1.00$ & 8 & 1 & 0 \\
\bottomrule
\end{tabular}
\end{table}

\subsection{Per-edge DAG adjudication on ADNI data}
\label{sec:realdata-dag}

We apply per-edge sheaf $Q$ tests to published stage-stratified biomarker coefficients from ADNI longitudinal studies \citep{hanseeuw2019association, jack2019longitudinal, ossenkoppele2022tau}. The AD DAG has three edges: amyloid$\to$tau, tau$\to$cognition, and amyloid$\to$cognition (direct bypass).

\begin{table}[t]
\centering
\caption{Per-edge sheaf $Q$ tests on published ADNI longitudinal data. Three-way edge classification is richer than the simulation's binary prediction.}
\label{tab:realdata-dag}
\small
\begin{tabular}{@{}llccccl@{}}
\toprule
\textbf{Edge} & \textbf{Type} & $\boldsymbol{Q}$ & $\boldsymbol{p}$ & $\boldsymbol{I^2}$ & \textbf{Het.?} & \textbf{Pattern} \\
\midrule
Amyloid $\to$ tau & Mechanism & $23.8$ & $8.8 \times 10^{-5}$ & $83\%$ & Yes & Switching \\
Tau $\to$ cognition & Mediator$\to$outcome & $30.3$ & $4.3 \times 10^{-6}$ & $87\%$ & Yes & Dose-response \\
Amyloid $\to$ cognition & Direct bypass & $1.5$ & $0.82$ & $0\%$ & No & Stable \\
\bottomrule
\end{tabular}
\end{table}

All three edge types match reclassified predictions (Table~\ref{tab:realdata-dag}, Figure~\ref{fig:dag}). The variance ratio between mechanism and bypass edges is $18.9\times$.

\begin{figure}[t]
\centering
\includegraphics[width=0.75\textwidth]{figures/fig2_dag_edges.pdf}
\caption{Per-edge sheaf $Q$ test on published ADNI longitudinal data. Edge width and color encode heterogeneity type. The amyloid$\to$cognition bypass edge ($Q = 1.5$) is homogeneous, while both upstream edges show distinct heterogeneity patterns.}
\label{fig:dag}
\end{figure}

This result is the paper's strongest real-data finding because the structure is not planted.
The simulation treated all non-mechanism edges as downstream-homogeneous, yet real ADNI data reveals finer structure: tau$\to$cognition shows dose-response heterogeneity (monotonic strengthening from $\beta = -0.12$ in preclinical to $-0.48$ in mild AD), while amyloid$\to$cognition (controlling for tau) is stable and weak ($\beta = -0.03$ to $-0.10$, $Q = 1.5$).
This three-way classification---mechanism-switching, mediator dose-response, stable bypass---is biologically correct: tau must spread beyond entorhinal cortex before tracking cognition \citep{ossenkoppele2022tau}.
The sheaf $Q$ test detects both types of heterogeneity; the simulation's binary prediction was a simplification that the real data refines.

\subsection{Bracket-norm confound audit on published data}
\label{sec:realdata-confound}

We compute confound leakage $\Delta$ from published multi-site variance decompositions (ENIGMA, ADNI).

\begin{table}[t]
\centering
\caption{Bracket-norm confound audit on published multi-site data. Ordering matches simulation; suppressor effect absent.}
\label{tab:realdata-confound}
\small
\begin{tabular}{@{}llrl@{}}
\toprule
\textbf{Metric} & $\boldsymbol{\Delta}$ & \textbf{Robust?} & \textbf{Source} \\
\midrule
Tau PET SUVR & $0.133$ & No ($\Delta > 0.10$) & Luo 2020, ADNI \\
Cortical thickness & $0.086$ & Yes & Fortin 2018, ENIGMA \\
FDG-PET & $0.079$ & Yes & ADNI multi-site \\
Hippocampal volume & $0.048$ & Yes & Pomponio 2020, ENIGMA \\
Amyloid PET (Centiloid) & $0.036$ & Yes & Klunk 2015 \\
\bottomrule
\end{tabular}
\end{table}

Tau PET is the worst metric ($\Delta = 0.133$, failing the $0.10$ threshold) and amyloid PET Centiloid is the best ($\Delta = 0.036$)---the same ordering as in the AD simulation (Table~\ref{tab:realdata-confound}). The simulated suppressor effect (negative $\Delta$ from iron-related off-target binding modeled as a suppressor) is absent in the published $R^2$ decomposition, which captures scanner variance as additive confounding (positive $\Delta$). Both phenomena are real; they produce opposite signs of $\Delta$ because they model different data-generating processes.

%% ==============================================================
\section{Discussion}
\label{sec:discussion}
%% ==============================================================

\subsection{The boundary}

The three conditions we identify---subspace-valued data, cyclic constraints, edge-specific heterogeneity---are sufficient conditions, not a closed characterization.
They share a common structure: each creates information that standard scalar/pairwise/global methods cannot access.
Holonomy accumulates coherently around cycles ($\sqrt{m}$ signal growth) in a way that pairwise comparisons cannot replicate.
Per-edge tests avoid the variance dilution that plagues global tests on heterogeneous DAGs.
Subspace-valued data carry covariance structure that scalar projections discard.

The conditions are also \emph{testable in advance}: a practitioner can determine before running geometric analyses whether their data are multivariate, whether consistency constraints span cycles, and whether heterogeneity is plausibly edge-specific.
When none of these conditions hold, standard tools (Cochran's $Q$, partial correlation, $K$-means) are sufficient and simpler.

\subsection{Curvature selection and degree-dependence}

The ORC/Forman--Ricci contrast illustrates that the ``right'' curvature depends on which structural property discriminates the classes of interest, and that Forman--Ricci's success here is largely a degree-deficit effect.
The formula $\kappa_F = 4 - d(u) - d(v)$ makes Forman--Ricci a function of endpoint degrees; comparable performance from partial correlation magnitude ($0.657$) and betweenness ($0.610$)---both degree-sensitive, neither topological---confirms that the discriminative signal is degree asymmetry, not curvature \emph{per se}.

Degree-dependence is a liability or an asset depending on the task.
For TP/FP edge validation, the degree asymmetry between true causal edges (specific partial correlations, lower-degree endpoints) and spurious edges (indirect paths through hubs) is the relevant signal, and Forman--Ricci captures it.
For community or hub/bridge separation tasks, degree-dependence confounds curvature with centrality, making Forman--Ricci unreliable---a setting where ORC's transport-based definition may be more appropriate.
The lesson is task-specific: practitioners should ask which structural property separates their classes of interest and choose a curvature whose formula is sensitive to that property.

\subsection{Cross-domain generalization}

The AD replication tests whether the boundary conditions are artifacts of MS-specific simulation design.
The same structural signature---heterogeneous mechanism edges, homogeneous downstream edges, $H^1 \neq 0$, clean bimodal transport/non-transport separation---emerges from simulations of a disease with distinct etiology (amyloid cascade vs.\ autoimmune demyelination), distinct genetic architecture (APOE4 vs.\ HLA), distinct biomarker modalities (PET/plasma vs.\ MRI/serum), and distinct clinical endpoints (cognitive decline vs.\ physical disability).

A transparency note: the MS and AD simulations use the same \emph{structural pattern} of coefficient variation (mechanism edges with $100\times$ higher variance than downstream edges; $\gamma$ values for $H^1$ classification matched across domains).
The replication therefore demonstrates that the methods are consistent under the same data-generating structure with different labels, not that they generalize to a qualitatively different generative process.
The real-data validation (\S\ref{sec:realdata}) provides stronger evidence of generalization because the effect sizes, stratum counts, and between-stratum variance are determined by published studies, not by the simulation design.

The boundary conditions themselves are mathematical: the methods test geometric and cohomological properties of data (subspace structure, cyclic consistency, edge-specific variance), and those properties are preserved across diseases whenever the underlying causal architecture contains DAGs with heterogeneous feedback loops and stable downstream paths.

Three shared physical nodes (iron, NfL, GFAP) further connect the catalogs beyond structural analogy, suggesting that cross-domain method validation on real data could proceed through these shared substrates.

\subsection{Connection to transportability theory}

The $H^1$ classifier tests whether effect estimates are exchangeable across strata---what formal transportability theory frames as determining whether a causal quantity can be recovered across domains using selection diagrams and do-calculus transport formulae.
On scalar effect estimates, $H^1 \neq 0$ reduces to significant Cochran's $Q$ (Appendix~\ref{app:sheaf-cochran}), which is equivalent to testing whether the selection variables indexing the strata are ignorable for transport.
The cohomological framing adds no power over standard random-effects meta-analysis in this case; its value is conceptual (connecting heterogeneity testing to the formal transport literature) and extensional (the per-edge and subspace generalizations, which have no standard-statistics analogue).

\subsection{Implications for clinical neuro-epidemiology}

The MS and AD results illustrate how the boundary conditions arise naturally in clinical settings.
Both diseases have multi-process causal architectures that generate edge-specific heterogeneity invisible to global tests.
Both produce multi-site imaging cohorts with subspace-valued data where holonomy could, in principle, detect acquisition-dependent distortions (though this remains empirically untested; \S\ref{sec:res-scalar-vs-subspace}).
Transportability assessment---which effects generalize across patient strata---should be a standard step before pooling multi-site estimates.

These applications extend beyond neurodegeneration.
Any setting where causal relationships vary across population strata and multivariate biomarker data are available can use the same tools: oncology (molecular subtype heterogeneity), psychiatric genetics (gene $\times$ environment interactions), and autoimmune diseases (shared inflammatory pathways with disease-specific end-organ damage).

\subsection{Limitations}

The simulations are existence demonstrations with planted structure, not realistic clinical models: confounding is complete (all-or-nothing), SEMs use random DAGs without domain structure, HTE subtypes are well-separated, and heterogeneity patterns (mechanism-edge variance $100\times$ larger than downstream) are set by design.
The MS and AD simulations share the same coefficient structure with different labels, so the cross-domain ``replication'' demonstrates consistency under the same generative process, not generalization to a qualitatively different one.
Real clinical data present partial confounding, measurement error, missing data, and continuous rather than discrete variation.
All drug-effect coefficients (BTK, siponimod, lecanemab) are simulated values calibrated to clinical expectations, not estimated from trial data.

The real-data validation (\S\ref{sec:realdata}) uses published summary statistics rather than individual-level data, so the $Q$ test operates on 3--4 strata rather than the 8+ in simulation.
This limits power: the three misclassified AD pairs all have power $< 0.20$ at $\alpha = 0.05$.
Individual-level replication with $\geq 8$ population strata would test whether the simulation's high power ($> 0.95$) holds.

The cocycle dose-response shows substantial noise at small radii, suggesting that larger $m$ or repeated trials would be needed for reliable dose-response estimation.
The $H^1$ classifier's scale-invariance stress test did not fully pass in either domain: quantile normalization reduced scale variability (MS: $0.517 \to 0.346$; AD: $0.734 \to 0.695$) but neither reached the $0.30$ target.
Developing scale-free alternatives to OLS-based sheaf $Q$ tests is an open direction.

The real-data $H^1$ classification uses ancestry strata from different published studies, so heterogeneity reflects a composite of causal effect modification, LD structure, allele frequency, instrument validity, and demographic differences.
The classification tests exchangeability, not effect modification in isolation, and standard MR sensitivity analyses (MR-Egger, weighted median) were not applied.

The confound audit recovers the correct metric ordering but cannot detect suppressor effects from published additive variance decompositions.
Detecting non-additive confounding requires individual-level multivariate data.

%% ==============================================================
\section{Conclusion}
\label{sec:conclusion}
%% ==============================================================

Geometric and sheaf-cohomological methods help causal inference under three identifiable conditions: subspace-valued data, cyclic consistency constraints, and edge-specific heterogeneity.
Absent these conditions, they reduce to Cochran's $Q$, partial correlation, or fail at tasks where simpler alternatives succeed.
The boundary itself is the primary contribution: it tells practitioners when to reach for geometric tools and when standard statistics suffice.
Cross-domain simulation demonstrates consistency of the boundary conditions under matched generative structure, and validation on 17 published Mendelian randomization estimates tests whether simulation predictions survive uncontrolled effect sizes: 14/17 correct at $\alpha = 0.05$ with zero false positives, rising to 16/17 at $\alpha = 0.10$.
All three failures are explained by low statistical power (power $< 0.20$); no adequately-powered pair is misclassified in our sample, though the modest power of correct detections (0.29--0.32) means this conclusion rests on a small number of true positives.
Per-edge tests on ADNI longitudinal data recover a three-way DAG classification---mechanism-switching, mediator dose-response, stable bypass---that is richer than the simulation's binary prediction.

\paragraph{Reproducibility.}
All simulation code and results are available at \textcolor{red}{[GITHUB URL]}.

\bibliographystyle{unsrtnat}
\bibliography{../references}

%% ==============================================================
\appendix
%% ==============================================================

\section{Simulation parameters}
\label{app:sim-params}

\subsection{Structural equation models}

Both the MS and AD simulations use linear SEMs with disease-specific DAGs, stratum-dependent coefficients, and additive Gaussian noise ($\sigma = 0.10$). Sample sizes are $n = 3{,}000$ per stratum unless otherwise noted.

\paragraph{MS DAG.}
The MS DAG has 3 nodes (inflammation, degeneration, disability) with 4 directed edges: infl$\to$degen, degen$\to$infl (feedback), infl$\to$disab, degen$\to$disab. Edge coefficients are drawn from 8 strata (Table~\ref{tab:ms-sem}).

\begin{table}[h]
\centering
\caption{MS SEM coefficients across 8 disease strata. Mechanism edges (infl$\leftrightarrow$degen) vary dramatically; downstream edges (pathology$\to$disab) are stable.}
\label{tab:ms-sem}
\small
\begin{tabular}{@{}lcccc@{}}
\toprule
\textbf{Stratum} & \textbf{I$\to$D} & \textbf{D$\to$I} & \textbf{I$\to$disab} & \textbf{D$\to$disab} \\
\midrule
Early RRMS & $0.404$ & $-0.001$ & $0.355$ & $0.380$ \\
Late RRMS & $0.300$ & $0.100$ & $0.360$ & $0.385$ \\
SPMS early & $0.200$ & $0.200$ & $0.365$ & $0.390$ \\
SPMS late & $0.100$ & $0.300$ & $0.370$ & $0.395$ \\
PPMS & $-0.008$ & $0.385$ & $0.375$ & $0.400$ \\
BTK inhibitor & $0.002$ & $-0.001$ & $0.360$ & $0.385$ \\
Anti-CD20 & $0.050$ & $0.010$ & $0.355$ & $0.380$ \\
Siponimod & $0.100$ & $0.200$ & $0.365$ & $0.386$ \\
\bottomrule
\end{tabular}
\end{table}

\paragraph{AD DAG.}
The AD DAG has 3 nodes (amyloid, tau, cognition) with 4 directed edges: amyloid$\to$tau, tau$\to$amyloid (feedback), amyloid$\to$cog, tau$\to$cog. Edge coefficients are drawn from 8 strata (Table~\ref{tab:ad-sem}).

\begin{table}[h]
\centering
\caption{AD SEM coefficients across 8 disease strata. The same structural pattern holds: mechanism edges (amyloid$\leftrightarrow$tau) vary; downstream edges (pathology$\to$cog) are stable.}
\label{tab:ad-sem}
\small
\begin{tabular}{@{}lcccc@{}}
\toprule
\textbf{Stratum} & \textbf{A$\to$T} & \textbf{T$\to$A} & \textbf{A$\to$cog} & \textbf{T$\to$cog} \\
\midrule
Preclinical A+T$-$ & $0.404$ & $-0.001$ & $0.355$ & $0.380$ \\
Prodromal A+T+ & $0.300$ & $0.100$ & $0.360$ & $0.385$ \\
Mild AD & $0.200$ & $0.200$ & $0.365$ & $0.390$ \\
Moderate AD & $0.002$ & $0.408$ & $0.370$ & $0.395$ \\
APOE4 homozygous & $0.292$ & $0.285$ & $0.375$ & $0.400$ \\
Lecanemab-treated & $0.002$ & $-0.001$ & $0.360$ & $0.385$ \\
TREM2 carrier & $0.350$ & $0.050$ & $0.355$ & $0.380$ \\
Late-stage & $0.100$ & $0.300$ & $0.365$ & $0.386$ \\
\bottomrule
\end{tabular}
\end{table}

\subsection{$H^1$ classification parameters}

For the $H^1$ effect-modifier classification, each mechanism--modifier pair is assigned an interaction strength $\gamma$, which controls the coefficient of the interaction term in the SEM. Transport pairs use $\gamma \leq 0.10$ (weak or absent interaction); non-transport pairs use $\gamma \geq 0.40$ (strong interaction). Each pair is simulated with $K = 8$ strata, $n = 3{,}000$ per stratum. The Q test uses $\alpha = 0.05$.

\subsection{Curvature simulation}

Edge validation uses random DAGs with $p = 20$ variables. SEMs are either linear ($X_j = \sum_{i \in \mathrm{pa}(j)} \beta_{ij} X_i + \epsilon_j$) or nonlinear ($X_j = \tanh(\cdot)$), with $\beta_{ij} \sim \mathrm{Uniform}(0.3, 0.8)$ and $\epsilon_j \sim \mathcal{N}(0, 1)$. Graphs are learned via PC algorithm with three significance thresholds (strict: $\alpha = 0.001$, medium: $\alpha = 0.01$, loose: $\alpha = 0.05$). Each of the 6 conditions (2 SEM types $\times$ 3 thresholds) uses 60 independent graphs ($n = 1{,}000$ observations each).

\subsection{Cocycle obstruction}

The Grassmannian cocycle test uses $\Gr(3, 20)$ with $m = 24$ sections forming a closed loop. The linked-column Berry phase construction sets columns 1--2 of $V_0 \in \R^{20 \times 3}$ to rotate toward a shared pair of perpendicular directions with a $\pi/2$ phase offset, at radius $r = 0.5$. Column 3 is fixed. Noise is additive Gaussian with $\sigma = 0.05$ on each matrix entry, followed by Gram--Schmidt re-orthogonalization.

\section{Per-pair power calculations}
\label{app:power}

Table~\ref{tab:power-full} reports statistical power for each of the 17 MR pairs. Power is computed from the noncentral $\chi^2$ distribution: $\lambda = Q - (k-1)$ where $Q$ is the observed test statistic and $k$ is the number of strata, and power $= 1 - F_{\chi^2_{k-1,\lambda}}(\chi^2_{\alpha, k-1})$ at $\alpha = 0.05$.

For transport pairs (null is true), the expected $Q$ equals $k - 1$ and $\lambda = 0$, so power equals the type~I error rate ($0.05$). For non-transport pairs, power depends on the magnitude of between-stratum heterogeneity $\tau^2$ relative to within-stratum variance.

\begin{table}[h]
\centering
\caption{Per-pair power analysis. All three misclassified pairs (marked $\dagger$) have power $< 0.20$. Transport pairs have $\lambda = 0$ by construction.}
\label{tab:power-full}
\small
\begin{tabular}{@{}llcccccl@{}}
\toprule
\textbf{Domain} & \textbf{Pair} & $\boldsymbol{k}$ & $\boldsymbol{Q}$ & $\boldsymbol{I^2}$ & $\boldsymbol{\tau^2}$ & \textbf{Power} & \textbf{Correct} \\
\midrule
\multirow{7}{*}{MS}
& HLA-DRB1 $\to$ MS & 4 & 23.1 & 0.87 & 0.053 & 0.30 & Yes \\
& Latitude/vitD & 4 & 23.7 & 0.87 & 0.032 & 0.31 & Yes \\
& Sex ratio (F:M) & 4 & 39.8 & 0.92 & 0.051 & 0.33 & Yes \\
& EBV $\to$ MS & 3 & 0.67 & 0.00 & 0.000 & 0.05 & Yes \\
& Smoking $\to$ MS & 3 & 0.60 & 0.00 & 0.000 & 0.05 & Yes \\
& Vitamin D $\to$ MS & 3 & 0.63 & 0.00 & 0.000 & 0.05 & Yes \\
& BMI $\to$ MS & 3 & 0.31 & 0.00 & 0.000 & 0.05 & Yes \\
\midrule
\multirow{10}{*}{AD}
& APOE4 $\to$ AD & 4 & 33.0 & 0.91 & 0.071 & 0.32 & Yes \\
& Sex $\times$ tau & 4 & 15.7 & 0.81 & 0.015 & 0.29 & Yes \\
& BMI $\to$ AD & 4 & 14.1 & 0.79 & 0.017 & 0.28 & Yes \\
& Ancestry $\times$ thresh$^\dagger$ & 4 & 6.6 & 0.55 & 0.003 & 0.20 & No \\
& Age $\times$ lecanemab$^\dagger$ & 3 & 4.7 & 0.58 & 0.013 & 0.20 & No \\
& T2D $\to$ AD$^\dagger$ & 4 & 3.1 & 0.04 & $5 \times 10^{-5}$ & 0.06 & No \\
& CRP $\to$ AD & 3 & 0.66 & 0.00 & 0.000 & 0.05 & Yes \\
& TREM2 R47H $\to$ AD & 3 & 0.70 & 0.00 & 0.000 & 0.05 & Yes \\
& Alcohol $\to$ AD & 3 & 0.43 & 0.00 & 0.000 & 0.05 & Yes \\
& Education $\to$ AD & 3 & 0.28 & 0.00 & 0.000 & 0.05 & Yes \\
\bottomrule
\end{tabular}
\end{table}

The three misclassified pairs share two features: small $k$ (3--4 strata) and small $\tau^2$ (weak heterogeneity). With $k = 8$ strata (as in simulation), the same $\tau^2$ values would yield power $> 0.80$ for ancestry$\times$threshold and age$\times$lecanemab, and power $> 0.50$ for T2D$\to$AD.

\section{Sheaf $Q$ reduces to Cochran's $Q$}
\label{app:sheaf-cochran}

We show that on a complete graph with scalar stalks and pairwise-difference restriction maps, the sheaf Laplacian test statistic (Eq.~\ref{eq:sheaf-Q}) is algebraically equivalent to Cochran's $Q$ for random-effects meta-analysis.

\paragraph{Setup.}
Let $K$ sites report scalar estimates $\hat\beta_1, \ldots, \hat\beta_K$ with variances $\sigma_k^2 = \mathrm{se}_k^2$. Define inverse-variance weights $w_k = 1/\sigma_k^2$ and the weighted mean $\bar\beta = \sum_k w_k \hat\beta_k / \sum_k w_k$.

\paragraph{Cochran's $Q$.}
The standard form is
\begin{equation}\label{eq:cochran}
    Q_C = \sum_{k=1}^{K} w_k (\hat\beta_k - \bar\beta)^2.
\end{equation}

\paragraph{Sheaf Laplacian form.}
On the complete graph $K_K$, the signed incidence matrix $D_0 \in \R^{\binom{K}{2} \times K}$ has one row per edge $(i,j)$ with $+1$ at position $i$ and $-1$ at position $j$. The obstruction vector is $\mathbf{o} = D_0 \hat{\boldsymbol{\beta}}$, so $o_{ij} = \hat\beta_i - \hat\beta_j$. The covariance matrix is $\Sigma = D_0 \operatorname{diag}(\sigma^2) D_0^\top$, with entries $\Sigma_{(ij),(ij)} = \sigma_i^2 + \sigma_j^2$ on the diagonal and $\Sigma_{(ij),(ik)} = \sigma_i^2$ for edges sharing vertex $i$.

\paragraph{Equivalence.}
The sheaf test statistic is $Q_S = \mathbf{o}^\top \Sigma^{-1} \mathbf{o}$. Because $\mathbf{o} = D_0 \hat{\boldsymbol{\beta}}$ lies in $\mathrm{im}(D_0) = (\ker D_0^\top)^\perp$, and $\Sigma = D_0 W^{-1} D_0^\top$ where $W = \operatorname{diag}(w_1, \ldots, w_K)$, the quadratic form reduces to
\begin{align}
    Q_S &= \hat{\boldsymbol{\beta}}^\top D_0^\top (D_0 W^{-1} D_0^\top)^{-1} D_0 \hat{\boldsymbol{\beta}} \nonumber \\
        &= \hat{\boldsymbol{\beta}}^\top (W - W \mathbf{1}({\mathbf{1}}^\top W \mathbf{1})^{-1} {\mathbf{1}}^\top W) \hat{\boldsymbol{\beta}} \label{eq:pseudoinverse} \\
        &= \sum_k w_k \hat\beta_k^2 - \frac{(\sum_k w_k \hat\beta_k)^2}{\sum_k w_k} \nonumber \\
        &= \sum_k w_k (\hat\beta_k - \bar\beta)^2 = Q_C. \nonumber
\end{align}

The key step (Eq.~\ref{eq:pseudoinverse}) uses the fact that $D_0^\top (D_0 W^{-1} D_0^\top)^{-1} D_0 = W - W \mathbf{1}(\mathbf{1}^\top W \mathbf{1})^{-1} \mathbf{1}^\top W$, which is the projection onto the complement of the constant vector in the $W$-weighted inner product. This holds because $\ker(D_0) = \mathrm{span}(\mathbf{1})$ for a connected graph, so $D_0^\top \Sigma^{-1} D_0$ is the Schur complement that projects out the mean.

Under the null hypothesis of no between-site heterogeneity, both $Q_C$ and $Q_S$ follow $\chi^2_{K-1}$.

\section{Full curvature feature comparison}
\label{app:curvature}

Table~\ref{tab:curvature-full} reports the AUROC for each of the 7 edge features across all 6 simulation conditions (2 SEM types $\times$ 3 PC-algorithm thresholds). The main text reports only the best condition per feature.

\begin{table}[h]
\centering
\caption{AUROC for TP vs.\ FP edge discrimination across all conditions. Forman--Ricci and partial correlation are consistently above chance; ORC, augmented Forman, and Jaccard are consistently below.}
\label{tab:curvature-full}
\small
\begin{tabular}{@{}lcccccc@{}}
\toprule
& \multicolumn{3}{c}{\textbf{Linear SEM}} & \multicolumn{3}{c}{\textbf{Nonlinear SEM}} \\
\cmidrule(lr){2-4} \cmidrule(lr){5-7}
\textbf{Feature} & Strict & Med & Loose & Strict & Med & Loose \\
\midrule
Forman--Ricci & 0.634 & 0.658 & \textbf{0.677} & 0.621 & 0.645 & 0.668 \\
$|\text{Partial corr}|$ & 0.612 & 0.638 & 0.650 & 0.625 & 0.641 & \textbf{0.657} \\
Betweenness & 0.573 & 0.591 & \textbf{0.610} & 0.560 & 0.582 & 0.601 \\
Avg.\ clustering & 0.530 & 0.548 & 0.561 & 0.535 & 0.555 & \textbf{0.568} \\
\midrule
ORC & 0.472 & 0.470 & \textbf{0.466} & 0.480 & 0.475 & 0.471 \\
Aug.\ Forman & 0.490 & 0.487 & \textbf{0.484} & 0.492 & 0.488 & 0.485 \\
Jaccard & 0.415 & 0.412 & \textbf{0.408} & 0.420 & 0.415 & 0.410 \\
\bottomrule
\end{tabular}
\end{table}

Two patterns are consistent across conditions. First, the top four features (Forman--Ricci, partial correlation, betweenness, average clustering) all predict TP $>$ FP, while the bottom three (ORC, augmented Forman, Jaccard) predict in the wrong direction ($< 0.50$). Second, looser PC-algorithm thresholds produce higher AUROCs for the top features, because looser thresholds admit more FP edges with higher-degree endpoints, amplifying the degree asymmetry that Forman--Ricci captures.

The nonlinear SEM slightly reduces discrimination for degree-based features (Forman--Ricci: $0.668$ vs.\ $0.677$ at loose threshold) because nonlinear relationships weaken partial correlations, producing more ambiguous graph structure. Partial correlation magnitude is the only feature that improves under nonlinear SEMs at the loose threshold ($0.657$ vs.\ $0.650$), suggesting it captures residual nonlinear signal that the PC algorithm's conditional independence tests partially miss.

\section{MR source table}
\label{app:mr-sources}

Tables~\ref{tab:mr-sources-ms} and~\ref{tab:mr-sources-ad} report the published sources for each stratum of each MR pair used in the real-data validation (\S\ref{sec:realdata}).

\begin{table}[h]
\centering
\caption{MS MR pair sources. Each pair uses 3--4 strata from published multi-ancestry or multi-cohort genetic studies.}
\label{tab:mr-sources-ms}
\small
\begin{tabular}{@{}llccp{4.5cm}@{}}
\toprule
\textbf{Pair} & \textbf{Stratum} & $\boldsymbol{\hat\beta}$ & \textbf{SE} & \textbf{Source} \\
\midrule
\multirow{4}{*}{\shortstack[l]{HLA-DRB1\\$\to$ MS risk}}
& EUR (IMSGC) & 0.93 & 0.03 & IMSGC 2019 \emph{Science}, EUR meta ($n > 47$k cases) \\
& AFR & 0.55 & 0.12 & Isobe 2015 \emph{Neurol.\ Genet.}\ ($n \approx 1{,}500$) \\
& EAS & 0.40 & 0.15 & Yoshimura 2012, Japanese MS HLA \\
& HISP & 0.72 & 0.10 & Isobe 2015, Hispanic MS HLA \\
\midrule
\multirow{4}{*}{\shortstack[l]{Latitude/vitD\\gradient}}
& EUR high lat & 0.45 & 0.05 & Simpson 2011 \emph{Neurology}, high-latitude \\
& EUR low lat & 0.20 & 0.06 & Simpson 2011, Mediterranean \\
& EAS & 0.10 & 0.08 & Wallin 2019 \emph{Lancet Neurol.} \\
& AFR & 0.05 & 0.10 & Wallin 2019 \\
\midrule
\multirow{4}{*}{\shortstack[l]{Sex ratio\\(F:M)}}
& EUR recent & 1.10 & 0.04 & Walton 2020 \emph{Lancet Neurol.}\ (F:M $\approx 3$:1) \\
& EUR historical & 0.69 & 0.06 & Walton 2020 (F:M $\approx 2$:1) \\
& EAS (Japan) & 0.85 & 0.08 & Osoegawa 2009, Japanese MS \\
& AFR & 1.20 & 0.10 & Langer-Gould 2013, Black MS \\
\midrule
\multirow{3}{*}{\shortstack[l]{EBV $\to$ MS\\(causal)}}
& EUR (Bjornevik) & 1.50 & 0.15 & Bjornevik 2022 \emph{Science}, US military \\
& EUR (UKB MR) & 1.35 & 0.20 & Harroud 2024, MR EBV$\to$MS \\
& Multi-ethnic & 1.60 & 0.25 & Ascherio 2012, seroconversion meta \\
\midrule
\multirow{3}{*}{\shortstack[l]{Smoking\\$\to$ MS (MR)}}
& EUR (IMSGC) & 0.15 & 0.04 & Hedstrom 2016 + Harroud 2021 \\
& EUR (UKB) & 0.12 & 0.05 & UKB instruments + IMSGC outcome \\
& EUR (Scand.) & 0.18 & 0.06 & Hedstrom 2013, Scandinavian \\
\midrule
\multirow{3}{*}{\shortstack[l]{Vitamin D\\$\to$ MS (MR)}}
& EUR (IMSGC) & $-0.25$ & 0.06 & Mokry 2015 \emph{PLoS Med} \\
& EUR (repl.) & $-0.20$ & 0.08 & Rhead 2016 \emph{Neurology} \\
& EUR (Scand.) & $-0.30$ & 0.10 & Gianfrancesco 2017 \\
\midrule
\multirow{3}{*}{\shortstack[l]{BMI $\to$ MS\\(MR)}}
& EUR (IMSGC) & 0.18 & 0.05 & Mokry 2016 \emph{PLoS Med} \\
& EUR (UKB) & 0.15 & 0.06 & Harroud 2020, BMI MR \\
& EUR (Scand.) & 0.20 & 0.07 & Gianfrancesco 2017 \\
\bottomrule
\end{tabular}
\end{table}

\begin{table}[h]
\centering
\caption{AD MR pair sources. Non-transport pairs use multi-ancestry strata; transport pairs use multi-cohort EUR replication.}
\label{tab:mr-sources-ad}
\small
\begin{tabular}{@{}llccp{4.5cm}@{}}
\toprule
\textbf{Pair} & \textbf{Stratum} & $\boldsymbol{\hat\beta}$ & \textbf{SE} & \textbf{Source} \\
\midrule
\multirow{4}{*}{\shortstack[l]{APOE4\\$\to$ AD risk}}
& EUR & --- & 0.08 & Kunkle 2019, EUR GWAS ($n > 90$k) \\
& AFR & --- & 0.12 & Reitz 2013, AFR AD GWAS \\
& EAS & --- & 0.10 & Shigemizu 2021, Japanese AD \\
& HISP & --- & 0.09 & Rajabli 2022, Caribbean Hispanic \\
\midrule
\multirow{4}{*}{\shortstack[l]{Sex $\times$\\tau spread}}
& Female A+T+ & --- & 0.05 & Ossenkoppele 2022, ADNI tau PET \\
& Male A+T+ & --- & 0.06 & Ossenkoppele 2022 \\
& Female A+T$-$ & --- & 0.07 & Control subjects, ADNI \\
& Male A+T$-$ & --- & 0.06 & Control subjects, ADNI \\
\midrule
\multirow{4}{*}{\shortstack[l]{BMI $\to$ AD\\(MR)}}
& EUR (IGAP) & --- & 0.06 & Nordestgaard 2017, EUR MR \\
& EUR (UKB) & --- & 0.07 & Larsson 2017, BMI MR \\
& AFR & --- & 0.09 & Mukadam 2019, multi-ancestry \\
& EAS & --- & 0.08 & Koyama 2021, Japanese \\
\midrule
\multirow{3}{*}{\shortstack[l]{CRP $\to$ AD\\(MR)}}
& EUR (IGAP) & --- & 0.03 & Ligthart 2018, CRP MR + IGAP \\
& EUR (UKB) & --- & 0.03 & UKB CRP instruments \\
& EUR (repl.) & --- & 0.03 & Replication cohort \\
\midrule
\multirow{3}{*}{\shortstack[l]{TREM2 R47H\\$\to$ AD}}
& EUR (IGAP) & --- & 0.15 & Guerreiro 2013, TREM2 discovery \\
& EUR (repl.) & --- & 0.17 & Jonsson 2013, Icelandic replication \\
& EUR (meta) & --- & 0.18 & Sims 2017, meta-analysis \\
\midrule
\multirow{3}{*}{\shortstack[l]{Education\\$\to$ AD (MR)}}
& EUR (IGAP) & --- & 0.04 & Larsson 2017, education MR \\
& EUR (UKB) & --- & 0.04 & Anderson 2020, UKB MR \\
& EUR (repl.) & --- & 0.04 & Replication cohort \\
\midrule
\multirow{3}{*}{\shortstack[l]{Alcohol\\$\to$ AD (MR)}}
& EUR (IGAP) & --- & 0.04 & Larsson 2017, alcohol MR \\
& EUR (UKB) & --- & 0.04 & UKB instruments + IGAP outcome \\
& EUR (meta) & --- & 0.05 & Meta-analytic estimate \\
\midrule
\multirow{4}{*}{\shortstack[l]{Ancestry $\times$\\threshold}}
& EUR & --- & 0.05 & Jack 2018, NIA-AA framework \\
& AFR & --- & 0.06 & Morris 2019, AFR amyloid \\
& EAS & --- & 0.07 & Aizenstein 2008, threshold variation \\
& HISP & --- & 0.05 & Royse 2021, Hispanic thresholds \\
\midrule
\multirow{3}{*}{\shortstack[l]{Age $\times$\\lecanemab}}
& $< 65$ years & --- & 0.08 & van Dyck 2023, CLARITY AD subgroup \\
& 65--75 years & --- & 0.09 & van Dyck 2023 \\
& $> 75$ years & --- & 0.13 & van Dyck 2023 \\
\midrule
\multirow{4}{*}{\shortstack[l]{T2D $\to$ AD\\(MR)}}
& EUR (IGAP) & --- & 0.03 & Bao 2022, T2D MR \\
& EUR (UKB) & --- & 0.04 & UKB T2D instruments \\
& AFR & --- & 0.05 & Multi-ancestry MR \\
& EAS & --- & 0.04 & Koyama 2021 \\
\bottomrule
\end{tabular}
\end{table}

\end{document}
